\section{Probing Spatial Capabilities}
\label{sec:formal-framework}

We ask a simple question: \emph{what spatial information can be decoded from a
frozen visual representation, and where in the network does it emerge?}
Spatial understanding is not a single ability. A representation may expose
the geometry of the currently observed scene without maintaining the location
of an occluded object; likewise, a future representation may be predictable
because consecutive frames look similar, rather than because the model uses
the camera action. We therefore separate three complementary capabilities and
evaluate each with a task-specific readout while keeping the backbone frozen.

\begin{table*}[t]
  \centering
  \small
  \setlength{\tabcolsep}{5pt}
  \begin{tabular}{p{0.18\textwidth}p{0.27\textwidth}p{0.47\textwidth}}
    \toprule
    \textbf{Capability} & \textbf{Why it is needed} & \textbf{What the probe does} \\
    \midrule
    \textbf{Geometric perception}
      & Establishes whether the representation describes the observed 3D
        scene, before asking it to remember or predict anything.
      & Reads out metric depth, relative camera pose, and dense instance identity,
        and tests whether representations from two views reveal their visual
        overlap. \\
    \textbf{Egocentric persistence}
      & Distinguishes a persistent, camera-relative object state from a
        snapshot of what is currently visible.
      & Forms an appearance query when an object is visible, then localizes
        that object after it becomes hidden. The output is its azimuth,
        elevation, and distance in the camera frame at the end of the prefix;
        the probe receives neither camera pose nor the target-time mask or
        location. \\
    \textbf{Action-conditioned prediction}
      & Tests whether the representation supports spatial updates under
        motion, rather than only static geometry or temporal smoothness.
      & Predicts an independently encoded future representation from the
        observed prefix and camera motion. We test a one-step transition,
        sequential rollout at horizons $1,2,4$, and a counterfactual
        intervention in which different actions from the same state must lead
        to their corresponding future representations. \\
    \bottomrule
  \end{tabular}
  \caption{The three spatial capabilities evaluated in our benchmark. The
  names in this table are used consistently throughout the paper; repository
  implementation codes are not used as scientific task names.}
  \label{tab:spatial-capabilities}
\end{table*}

The three capabilities answer different questions and are not assumed to form
a score hierarchy. In particular, geometric perception does not imply object
persistence, and high future-feature retrieval does not by itself establish
the use of action. We therefore accompany the prediction tasks with matched
no-action, shuffled-action, and last-observation-copy controls, and the
persistence task with object-conditioned random-feature and unconditional
controls.

\subsection{Causal streaming evaluation}
\label{sec:streaming-evaluation}

All temporal claims use posed ScanNet++ trajectories
\citep{yeshwanth2023scannetpp}. For model $m$, extraction layer $\ell$, prefix
length $p$, and current time $t$, we compute
\begin{equation}
  \mathbf{Z}^{m,\ell}_{t,p}
  = E^{m}_{\ell}\!\left(\mathbf{I}_{t-p+1:t}\right),
  \qquad p\in\{8,12,16,24\}.
  \label{eq:causal-streaming-feature}
\end{equation}
Each prefix is forwarded independently and contains no frame after $t$; the
four prefix lengths are trained and evaluated as separate fixed-shape
settings. This design measures how a representation changes as more past
observations become available without giving the encoder future context.

The leakage-sensitive inputs are constructed separately. For egocentric
persistence, the same object is selected from the common visible history and
reused across the prefix sweep, while its target is expressed in the camera
coordinate system at the current prefix tail. For action-conditioned
prediction, every future target at $t+h$ is produced by an independent
target-only encoder forward,
\begin{equation}
  \mathbf{Z}^{m,\ell,\mathrm{tar}}_{t+h}
  = E^{m}_{\ell}\!\left(\mathbf{I}_{t+h}\right),
  \qquad h\in\{1,2,4\},
  \label{eq:isolated-target-feature}
\end{equation}
so it cannot contain prefix tokens or later frames. Stochastic backbones use
fixed, matched noise seeds within a comparison. Splits are scene-disjoint.

\subsection{Models and extraction layers}
\label{sec:model-layer-grid}

We probe representations from image encoders, video encoders, 3D encoders,
video generation/world models, and a multimodal language model. Table
\ref{tab:model-layer-grid} lists the extraction points in the experimental
matrix. Dense layer sweeps are used for the principal video backbones, while
single canonical representations serve as reference points for the remaining
models. Layer numbers are model-specific extraction indices and should not be
compared as absolute depths across architectures.

\begin{table*}[t]
  \centering
  \small
  \setlength{\tabcolsep}{5pt}
  \begin{tabular}{p{0.20\textwidth}p{0.25\textwidth}p{0.25\textwidth}p{0.22\textwidth}}
    \toprule
    \textbf{Model family} & \textbf{Model} & \textbf{Extraction layers} &
    \textbf{Extraction condition} \\
    \midrule
    Image encoder
      & DINOv2-Large
      & final hidden state
      & deterministic \\
    Video encoder
      & V-JEPA2 ViT-L
      & $\{0,5,11,17,23\}$
      & deterministic \\
    Video encoder
      & VideoMAE-Large
      & $\{0,5,11,17,23\}$
      & deterministic \\
    3D encoder
      & Fast3R ViT-L
      & $\{24\}$
      & deterministic \\
    Video generator
      & Wan2.1-T2V-1.3B
      & $\{0,5,10,15,20,25,29\}$
      & scheduler step $749$ \\
    Video generator
      & CogVideoX-5B-I2V
      & $\{0,5,10,15,20,25,30,35,40\}$
      & scheduler step $749$ \\
    World model
      & AetherV1
      & $\{1\}$
      & scheduler step $749$ \\
    Video generator
      & MiniMax-H3
      & $\{10,20,30,40,50\}$
      & flow coordinate $0.25$ \\
    Video generator
      & LingBot-Video MoE-30B-A3B
      & $\{10,20,30,40,48\}$
      & flow coordinate $0.25$ \\
    Multimodal language model
      & Qwen2.5-VL-7B-Instruct
      & visual-merger output
      & deterministic visual forward \\
    \bottomrule
  \end{tabular}
  \caption{Models and model-specific extraction layers in the current
  experimental matrix. For DINOv2, the cache label $0$ denotes the final
  hidden state; for Qwen2.5-VL, cache label $-1$ denotes the visual-merger
  output. The terminal sweep labels $50$ for MiniMax-H3 and $48$ for
  LingBot-Video map to their final blocks (indices $49$ and $47$,
  respectively).}
  \label{tab:model-layer-grid}
\end{table*}

Within each model and capability, the readout architecture and training
protocol are held fixed across layers and prefix lengths. This produces a
layer--context map of each capability without assuming that the last layer or
the longest prefix is necessarily best.
